\section{联合调度约束条件与可行域}
\label{sec:constraints}

\subsection{物理约束汇总}

联合模型的可行域由4.2--4.7发布的设备方程和公共接口共同组成。源侧满足
式~\eqref{eq:source-use-curt}和各设备状态/爬坡边界；电池满足
式~\eqref{eq:bess-state}--\eqref{eq:gfm-down-reserve}；制氢、储氢和可选氢转电满足
式~\eqref{eq:ely-module}--\eqref{eq:ely-h2p-exclusive}；算力满足4.6的任务、队列、容量、
PUE与SLA约束；产品输出满足式~\eqref{eq:cable-send}--\eqref{eq:marine-service}；所有
动作最终满足式~\eqref{eq:master-bus}。

设备可用率或事件状态统一进入容量上界：
\begin{equation}
0\le x_{j,t}\le A_{j,t}^{evt}\,\overline x_{j,t},
\qquad j\in\mathcal J^{asset},
\label{eq:availability-general}
\end{equation}
其中$A_{j,t}^{evt}\in[0,1]$由当前观测、故障、检修、台风预警/过境/恢复状态或保护逻辑
给出。关闭某条产品路径时，其专属功率、流量和收入变量均置零；共同源侧、BESS和内部
刚性负荷不因产品路径关闭而消失。

\subsection{逐小时E/H/C固定策略约束}

固定对比策略的E/H/C比例作用于每个实际发生可选产品分配的小时，不使用提前已知的全期
总供电量。定义
\begin{equation}
p_{E,t}=P_t^{cab,s},\qquad
p_{H,t}=P_t^{EL},\qquad
p_{C,t}=P_t^{comp,f},\qquad
p_{\Sigma,t}=p_{E,t}+p_{H,t}+p_{C,t}.
\label{eq:hourly-mix-def}
\end{equation}
基础算力是所有方案共有的刚性负荷，不进入$p_{C,t}$及比例分母。给定目标比例
$q_E+q_H+q_C=1$和允许偏差$\delta$，对$i\in\{E,H,C\}$逐时施加
\begin{equation}
\max(0,q_i-\delta)p_{\Sigma,t}
\le p_{i,t}\le
\min(1,q_i+\delta)p_{\Sigma,t}.
\label{eq:hourly-mix}
\end{equation}
若$q_i=0$，则相应产品路径功率为零。$p_{\Sigma,t}=0$时比例约束自然退化为零分配，不
要求为满足比例而强制消耗能源。$\delta$及任何最小有效分配门槛均为
\assumptiontag。式~\eqref{eq:hourly-mix}只定义运行电能分配，不代表设备装机比例、投资
比例或产品收入比例。

\subsection{服务等级、终端与事件约束}

内部关键负荷、刚性海洋用能和刚性算力分别设置未供变量，并通过目标函数中的高惩罚或
式~\eqref{eq:epsilon-constraint}的ENS上限约束。惩罚系数只能排序可避免的缺供，不能
替代真实容量；资源和储备不足时，模型应显式输出ENS，不得通过提高惩罚系数隐藏不可行性。

有限时域可按研究目的选择以下终端规则：
\begin{equation}
E_T^B=E_0^B\ \text{或}\ E_T^B\ge E^{B,target},\qquad
H_T=H_0\ \text{或}\ H_T\ge H^{target}.
\label{eq:terminal-rules}
\end{equation}
循环规则适用于独立典型周期；最低目标适用于检修或安全储备；因果逐时年度运行采用状态
连续传递并另设储备策略，不在每小时强制循环。不同终端规则的结果不可直接比较。

极端事件使用以下外生状态：
\begin{equation}
z_t^{evt}\in\{\mathrm{NORMAL},\mathrm{WARNING},
\mathrm{PASSAGE},\mathrm{RECOVERY}\}.
\label{eq:event-state}
\end{equation}
在预警状态，
策略可暂停可选外送、制氢交付和柔性算力以恢复SOC与安全库存；过境状态按设备保护停机、
海缆/物流/通信可用率和刚性服务优先级运行；恢复状态按实测可用率逐步解除限制。具体事件
时间、降额、储备目标和恢复斜率须来自气象海洋预报、OEM/EPC和应急预案，未取得前为
\assumptiontag。

\subsection{变量域与线性化}

连续变量包括功率、电量、产量、库存、任务完成量和ENS；整数变量包括电池充放状态、
电解槽在线模块数、设备启动、算力在线状态和可选氢转电状态。固定SEC、线性库存递推、
分段线性PUE/损耗代理与线性产品比例使运行问题可写成MILP。任何分段线性代理均须在
求解后用原物理函数复算，若复算越界则收紧分段或可行域后重求。

本节不再设置“容量规划与资本回收”决策。资产建设容量作为场景输入，生命周期成本在
4.8独立分账；未来若开展容量--运行联合优化，应另增规划层和年度/多年度代表场景，不能
把当前运行MILP结果改名为容量最优。
